% !TEX program = xelatex
% 二次型与主轴 · 独立改写片段（2026-11-11）
\documentclass[11pt]{ctexart}
%=====================================================================
% theme-violet.tex —— 讲义精讲 · 紫花主题
% 用法：\documentclass[11pt]{ctexart}
%       \input{assets/theme-violet}
%       \begin{document} ... \end{document}
% 只用 ctexart 自带能力 + xcolor/enumitem/fancyhdr/pifont/amsmath，
% 不依赖 tikz/tcolorbox，保证 MiKTeX 可直接编译。
%=====================================================================

\usepackage[a4paper,left=2.0cm,right=2.0cm,top=2.1cm,bottom=2.3cm,%
            headsep=12pt,footskip=22pt]{geometry}
\usepackage{amsmath,amssymb}
\usepackage{needspace}
\usepackage{xcolor}
\usepackage{array}
\usepackage{enumitem}
\usepackage{pifont}
\usepackage{fancyhdr}
\usepackage[hidelinks]{hyperref}

%---------------------------------------------------------------- 配色
\definecolor{violet}{RGB}{107,78,143}      % 主色：紫
\definecolor{violetdeep}{RGB}{63,44,92}    % 深紫：标题/强调
\definecolor{petal}{RGB}{183,159,208}      % 花瓣：装饰线、花饰
\definecolor{bandviolet}{RGB}{247,243,251} % 浅紫底：档案/提示框
\definecolor{bandcream}{RGB}{252,249,242}  % 米白底：易错/注意框
\definecolor{rulepurple}{RGB}{217,207,230} % 分隔线
\definecolor{inkgray}{RGB}{92,96,110}      % 次要文字

%---------------------------------------------------------------- 字体
% 有 Noto CJK 用 Noto，没有就退回 ctex 默认（Windows 下为宋体/黑体）
\IfFontExistsTF{Noto Serif CJK SC}{\setCJKmainfont{Noto Serif CJK SC}}{}
\IfFontExistsTF{Noto Sans CJK SC}{\setCJKsansfont{Noto Sans CJK SC}}{}

%---------------------------------------------------------------- 版式
\setlist{nosep,leftmargin=1.9em}
\linespread{1.28}
\setlength{\parskip}{3.5pt}
\setlength{\parindent}{0pt}
\setlength{\tabcolsep}{5pt}
\renewcommand{\arraystretch}{1.28}

%---------------------------------------------------------------- 花饰
% \ding{95} = 五瓣小花，\ding{96} = 星形花；都来自 pifont，必定存在
\newcommand{\orn}{\textcolor{petal}{\ding{95}}}          % 行内小标
\newcommand{\ornstar}{\textcolor{petal}{\ding{96}}}      % 星形花
\newcommand{\flowerrule}{%                                % 居中花饰分隔线
  \par\vspace{7pt}%
  \noindent\makebox[\linewidth]{\color{petal}%
    \rule{0.20\linewidth}{0.6pt}\hspace{8pt}{\small\ding{95}}\hspace{4pt}%
    \raisebox{1.1pt}{\tiny\ding{95}}\hspace{4pt}{\small\ding{95}}\hspace{8pt}%
    \rule{0.20\linewidth}{0.6pt}}%
  \par\vspace{7pt}}

%---------------------------------------------------------------- 标题与节头
\newcommand{\doctitle}[3]{%   #1 主标题  #2 副标题  #3 元信息（可用 \\ 分行）
  \begin{center}
    {\color{petal}\small \orn\enspace #2\enspace \ornstar}\\[7pt]
    {\LARGE\bfseries\color{violetdeep} #1}\\[10pt]
    {\color{petal}\rule{0.55\linewidth}{1pt}}\\[4pt]
    {\small\color{inkgray} #3}
  \end{center}\vspace{2pt}}

\newcommand{\hdr}[2]{%   #1 章节号（如 §0）  #2 标题
  \par\vspace{12pt}%
  \noindent\colorbox{violetdeep}{\color{white}\bfseries\small\;#1\;}\hspace{7pt}%
  {\color{violetdeep}\large\bfseries #2}\par\nobreak
  \vspace{3pt}{\color{petal}\hrule height 1.2pt}\vspace{6pt}}

\newcommand{\subhdr}[1]{%
  \par\vspace{6pt}%
  \noindent{\color{violet}\bfseries #1}\par\vspace{1pt}}

%---------------------------------------------------------------- 提示框
\newcommand{\callout}[1]{% 浅紫底：讲义档案、读法说明
  \par\medskip\noindent
  \fcolorbox{rulepurple}{bandviolet}{\parbox{\dimexpr\linewidth-2\fboxsep-2\fboxrule}{\vspace{2pt}#1\vspace{2pt}}}%
  \par\medskip}

\newcommand{\warnbox}[1]{% 米白底：易错、边界条件
  \par\medskip\noindent
  \fcolorbox{rulepurple}{bandcream}{\parbox{\dimexpr\linewidth-2\fboxsep-2\fboxrule}{\vspace{2pt}#1\vspace{2pt}}}%
  \par\medskip}

\newcommand{\gapitem}[2]{% #1 缺口类型（紫） #2 内容
  \par\vspace{2pt}\noindent
  \hangindent=1.4em\hangafter=1
  {\color{petal}\small\ding{96}}\enspace{\bfseries\color{violetdeep}#1}\enspace #2\par}

% 题干块（\qstem）：每道例题／作业的「详解」开头必须先复述题干与出处。
% 起因：2026-10-25 第五章讲义把七道作业的题面集中放在一页、详解另起一页，
% 读者翻到「第 4 题」时看不到条件，等于一道没有题干的题目。
% 规则：详解块内自带题干，读者不翻页就能读懂题在问什么。
\newcommand{\qstem}[1]{%
  \par\medskip\noindent
  \fcolorbox{violet}{bandviolet}{%
    \parbox{\dimexpr\linewidth-2\fboxsep-2\fboxrule}{%
      \vspace{2pt}%
      {\small\bfseries\color{violet}题干}\hspace{0.8em}#1%
      \vspace{2pt}}}%
  \par\medskip}

% v2: 保护题干与开头，长题可跨页，续页补上下文。
\newcommand{\qbreak}{\needspace{6\baselineskip}}
% 独立练习：ID、来源、题干、考点、答案标签（通常 ans:ID）。
\newcommand{\practice}[5]{%
  \qbreak\subhdr{练习 #1}\qstem{（#2）#3}%
  \textbf{考点}：#4\par
  详解见附录第~\pageref{#5}~页。\par\medskip}
% 附录开头：ID、来源、条件摘要；自动产生 ans:ID 标签。
\newcommand{\answerstart}[3]{%
  \qbreak\subhdr{#1\quad 详解}\label{ans:#1}\qstem{（#2）#3}}
\newcommand{\answercontinue}[2]{%
  \needspace{5\baselineskip}\subhdr{#1\quad 续解}%
  {\small\color{inkgray}本页接续：#2}\par\medskip}
% 延迟提供默认值，让旧讲义可在导言区保留自己的 \newcommand{\step}。
\AtBeginDocument{\providecommand{\step}[1]{\needspace{3\baselineskip}\par\medskip\textbf{#1}\par}}

% 前提旁注：只放在「需要理解的那一步」后面，给陈述 + 本题代入，不重证。
\newcommand{\aside}[1]{%
  \par\vspace{1pt}\noindent
  \hangindent=1.6em\hangafter=1
  {\small\color{inkgray}%
    \hspace*{0.6em}{\color{petal}$\triangleright$}\hspace{2pt}#1}%
  \par\vspace{2pt}}

%---------------------------------------------------------------- 表格
% 本机 MiKTeX 的 colortbl 与 array 版本不匹配（\rowcolor 会报错），
% 故表格统一用「上粗紫线 + 表头 + 中细线」的三线式，不依赖 colortbl。
\newcommand{\tabtop}{\par\vspace{3pt}{\color{violet}\hrule height 1.2pt}\vspace{1pt}}
\newcommand{\tabmid}{\par\vspace{1pt}{\color{petal}\hrule height 0.5pt}\vspace{1pt}}
\newcommand{\tabbot}{\par\vspace{3pt}{\color{rulepurple}\hrule height 0.5pt}\vspace{8pt}}
\newcommand{\thead}[1]{{\bfseries\color{violetdeep}#1}}

%---------------------------------------------------------------- 页眉页脚
\pagestyle{fancy}
\fancyhf{}
\renewcommand{\headrulewidth}{0.4pt}
\renewcommand{\headrule}{\hbox to\headwidth{\color{rulepurple}\leaders\hrule height \headrulewidth\hfill}}
\fancyhead[L]{\small\color{inkgray}\songti \doctag}
\fancyhead[R]{\small\color{inkgray}\songti \coursetag}
\fancyfoot[C]{\small\color{inkgray}\orn\enspace{\color{violet}\thepage}\enspace\orn}

% 每篇讲义在导言区设这两个值
\newcommand{\doctag}{}
\newcommand{\coursetag}{}
\newcommand{\settags}[2]{\renewcommand{\doctag}{#1}\renewcommand{\coursetag}{#2}}

%---------------------------------------------------------------- 数学记号
\newcommand{\vect}[1]{\boldsymbol{#1}}
\newcommand{\dd}{\,\mathrm{d}}
\newcommand{\ddx}{\frac{\mathrm{d}}{\mathrm{d}x}}
\addtolength{\jot}{2pt}



\usepackage{tikz}
\usetikzlibrary{arrows.meta,calc}
\setlength{\headheight}{22pt}
\linespread{1.24}
\raggedbottom
\settags{专题 · 二次型与主轴}{理论力学}

% 记号：矢量粗斜体，矩阵粗正体
\newcommand{\ten}[1]{\mathbf{#1}}
\newcommand{\unitd}{\hat{\ten{1}}}
\newcommand{\tr}{\operatorname{tr}}
\newcommand{\R}{\mathbb{R}}
\newcommand{\figcap}[1]{\par\vspace{2pt}{\small\color{inkgray}#1}\par}
\newcommand{\step}[1]{\par\smallskip\noindent{\color{violetdeep}\bfseries #1}\hspace{0.35em}}
\tikzset{
  ax/.style={inkgray,densely dotted,->},
  vv/.style={thick,->,>={Stealth[length=1.7mm]}},
  vviolet/.style={violet,thick,->,>={Stealth[length=1.7mm]}},
}

\let\orighdr\hdr
\renewcommand{\hdr}[2]{\par\needspace{10\baselineskip}\orighdr{#1}{#2}}
\let\origsubhdr\subhdr
\renewcommand{\subhdr}[1]{\par\needspace{3\baselineskip}\origsubhdr{#1}}

\setlength{\abovedisplayskip}{6pt plus 2pt minus 2pt}
\setlength{\belowdisplayskip}{6pt plus 2pt minus 2pt}
\setlength{\emergencystretch}{1.5em}

\begin{document}

\doctitle{二次型与主轴}
{从椭球、特征值到实对称矩阵的谱定理}
{专题讲义片段\quad 2026-11-11\quad 来源：TM2026-Topic8-P1 p9；RIT Whelan notes 10 pp.8--10}

\callout{\textbf{本片段做什么。}
本片段只处理一个数学问题：给定三维实对称矩阵 $\ten T$，二次型 $E(\vect a)=\vect a\cdot\ten T\vect a$ 在什么坐标轴下没有交叉项？这些轴为什么一定存在、为什么互相垂直？答案就是谱定理。它是后续惯量张量与应力张量的共同数学底座；本片段只在 §1.6 指出接口，不展开那些章节的构造与平衡方程。
\par\smallskip
\textbf{取材与边界。}
TM2026-Topic8-P1 p9 把二阶张量同时读成线性算子 $\vect b=\ten T\vect a$ 与二次型 $E=\vect a\ten T\vect a$；RIT Whelan notes 10 的 pp.8--10 给出“实对称张量可正交对角化”的结论、特征方程、特征值为实与不同特征向量正交的断言。原页没有给出奇次实多项式引理、$2\times2$ 判别式引理与正交补降维的证明；本片段 §1.4 的证明路线（先 $2\times2$ 引理，再三维实特征值，再正交补降维，最后拼基）是本片段自行补足的论证，不归给原页。
\par\smallskip
\textbf{读法。}
§0 只补后面真正会用到的工具：点积、长度、正交与投影，矩阵作用与转置，行列式与非零解判据，奇次实多项式有实根；特征值、特征向量与特征方程在 §1.2--§1.3 正式建立。§1 是主体；1.4 的证明较长，按“$2\times2$ 引理 $\to$ $3\times3$ 存在实特征值 $\to$ 正交补降维 $\to$ 拼基”读。P7 是独立练习，正文只给题干，详解在附录。}

\warnbox{\textbf{记号。}
矢量用粗斜体 $\vect a$，矩阵用粗正体 $\ten T$；$\ten T^{\mathsf T}$ 是转置，$\unitd$ 是单位矩阵，$\tr$ 是迹，$\vect a\cdot\vect b$ 是点积。三维指标从 $1$ 到 $3$，例如 $a_i$ 是 $\vect a$ 的第 $i$ 个分量。全文只讨论实数。}

\hdr{§0}{读前工具：只补本节要用的}

这一节不是线性代数复习课。它只建立后面反复使用的四件事：怎样用点积判断垂直、矩阵与转置怎样配合、行列式为什么能判定非零解、以及一个三次实系数多项式为什么一定有实根。每一项都直接服务于 §1 的某一步，末尾会说明服务的是哪一步。

\subhdr{0.1\quad 点积、长度、正交与正交归一基}

取三维实矢量
\[
\vect a=(a_1,a_2,a_3),\qquad \vect b=(b_1,b_2,b_3).
\]
它们的\textbf{点积}是一个数：
\[
\vect a\cdot\vect b=a_1b_1+a_2b_2+a_3b_3=\sum_{i=1}^{3}a_ib_i .
\]
点积满足交换律 $\vect a\cdot\vect b=\vect b\cdot\vect a$，并且对每个变量都是线性的。

\textbf{长度。}矢量自身的点积
\[
|\vect a|^2=\vect a\cdot\vect a=a_1^2+a_2^2+a_3^2
\]
是三个平方之和，所以 $|\vect a|^2\ge 0$，且只有 $\vect a=\vect 0$ 时才等于 $0$。定义
\[
|\vect a|=\sqrt{\vect a\cdot\vect a}.
\]
这就是通常的几何长度：在三个坐标轴互相垂直的坐标系里，连续用两次勾股定理，原点到点 $(a_1,a_2,a_3)$ 的距离正好是 $\sqrt{a_1^2+a_2^2+a_3^2}$。长度为 $1$ 的矢量叫\textbf{单位矢量}。

\textbf{点积与夹角。}设 $\vect a,\vect b$ 都非零，$\theta$ 是它们之间的夹角（把两矢量平移到同一起点后张开的角度）。在 $\vect a,\vect b$ 张成的平面里，第三边是 $\vect a-\vect b$。由平面几何的余弦定理（在 $\vect a,\vect b$ 张成的平面里作高，用勾股定理即可得到这条公式），
\[
|\vect a-\vect b|^2=|\vect a|^2+|\vect b|^2-2|\vect a||\vect b|\cos\theta .
\]
另一方面，用点积展开：
\[
|\vect a-\vect b|^2
=(\vect a-\vect b)\cdot(\vect a-\vect b)
=|\vect a|^2+|\vect b|^2-2\,\vect a\cdot\vect b .
\]
两式比较，得到
\[
\vect a\cdot\vect b=|\vect a||\vect b|\cos\theta .
\]
这就是点积的几何含义：它把两个长度和夹角合成为一个数。特别地，若 $\vect n$ 是单位矢量，则 $\vect a\cdot\vect n=|\vect a|\cos\theta$，其中 $\theta$ 是 $\vect a$ 与 $\vect n$ 的夹角。

\textbf{垂直判据。}若 $\vect a\cdot\vect b=0$ 且 $\vect a,\vect b$ 非零，则 $|\vect a||\vect b|\ne0$，所以 $\cos\theta=0$，即 $\theta=90^\circ$，两矢量垂直（正交）。反过来，若 $\vect a\perp\vect b$，则 $\cos\theta=0$，所以 $\vect a\cdot\vect b=0$。因此
\[
\vect a\cdot\vect b=0
\quad\Longleftrightarrow\quad
\vect a\text{ 与 }\vect b\text{ 正交}\qquad(\vect a,\vect b\ne\vect0).
\]
“点积为零”不是数值碰巧为零，它正好是夹角余弦为零，所以几何上就是直角。

\textbf{具体核对。}在平面里取 $\vect a=(3,0)$、$\vect b=(3,4)$。则
\[
|\vect a|=3,\qquad |\vect b|=5,\qquad \vect a\cdot\vect b=9 .
\]
由 $\vect a\cdot\vect b=|\vect a||\vect b|\cos\theta$ 得 $\cos\theta=9/(3\cdot5)=3/5$，所以 $\theta\approx53.1^\circ$；$\vect b$ 在 $\vect a$ 方向上的投影长是 $|\vect b|\cos\theta=5\cdot(3/5)=3$，正好是 $\vect b$ 的水平分量。再取 $\vect c=(0,4)$：$\vect a\cdot\vect c=0$，所以 $\vect a\perp\vect c$，夹角是 $90^\circ$。前面常用的例子 $(1,2,2)\cdot(2,-1,0)=2-2+0=0$ 用的是同一判据；长度 $3$ 与 $\sqrt5$ 只是核对。

\textbf{正交归一基。}一组两两正交的单位矢量 $\vect e_1,\vect e_2,\vect e_3$ 叫\textbf{正交归一基}。它满足
\[
\vect e_i\cdot\vect e_j=\delta_{ij}=
\begin{cases}
1,&i=j,\\
0,&i\ne j.
\end{cases}
\]
在正交归一基下，矢量的分量可以用点积直接读出来。若
\[
\vect a=a_1\vect e_1+a_2\vect e_2+a_3\vect e_3=\sum_{k=1}^{3}a_k\vect e_k,
\]
则两边与 $\vect e_j$ 作点积，利用 $\vect e_k\cdot\vect e_j=\delta_{kj}$，得到
\[
\vect a\cdot\vect e_j=\sum_{k=1}^{3}a_k\delta_{kj}=a_j .
\]
所以 $\vect a=\sum_j(\vect a\cdot\vect e_j)\vect e_j$。这不是新公式，而是“分量就是投影”的严格写法。

\textbf{沿单位方向的投影。}设 $\vect n$ 是单位矢量，$\vect a$ 任意。把 $\vect a$ 写成
\[
\vect a=(\vect a\cdot\vect n)\vect n+\vect r,
\qquad
\vect r=\vect a-(\vect a\cdot\vect n)\vect n .
\]
第一项沿 $\vect n$；第二项与 $\vect n$ 正交，因为
\[
\vect r\cdot\vect n
=\vect a\cdot\vect n-(\vect a\cdot\vect n)(\vect n\cdot\vect n)
=\vect a\cdot\vect n-\vect a\cdot\vect n=0 .
\]
所以 $\vect a\cdot\vect n$ 是 $\vect a$ 在 $\vect n$ 方向上的\textbf{带符号投影长度}：若 $\vect a\cdot\vect n>0$，投影与 $\vect n$ 同向；若小于 $0$，反向；它的绝对值才是投影的长度。投影矢量是 $(\vect a\cdot\vect n)\vect n$。这与前面的 $\vect a\cdot\vect n=|\vect a|\cos\theta$ 一致。这个“减去投影就得到正交补”的动作，§1.4 证明 $\ten T$ 保持正交补时会再用一次。

\textbf{小例。}取 $\vect a=(1,2,2)$，$\vect n=(1,1,1)/\sqrt3$。则
\[
\vect a\cdot\vect n=\frac{1+2+2}{\sqrt3}=\frac{5}{\sqrt3},
\qquad
(\vect a\cdot\vect n)\vect n=\frac53(1,1,1).
\]
于是
\[
\vect a-\frac53(1,1,1)=\left(-\frac23,\frac13,\frac13\right),
\]
它与 $\vect n$ 的点积为
\[
\frac{1}{\sqrt3}\left(-\frac23+\frac13+\frac13\right)=0 .
\]
这里 $\vect a\cdot\vect n=5/\sqrt3>0$，所以投影与 $\vect n$ 同向；剩余部分落在与 $\vect n$ 垂直的平面里。这就是“把 $\vect a$ 分成沿 $\vect n$ 的部分和垂直于 $\vect n$ 的部分”。

\subhdr{0.2\quad 矩阵、转置与对称}

一个 $3\times3$ 实矩阵 $\ten T=(T_{ij})$ 作用在矢量 $\vect x=(x_1,x_2,x_3)$ 上，得到一个新矢量 $\ten T\vect x$，其第 $i$ 个分量是
\[
(\ten T\vect x)_i=\sum_{j=1}^{3}T_{ij}x_j=T_{ij}x_j .
\]
这里用了求和约定：同一项里同一个指标出现两次，就对 $1,2,3$ 求和。矩阵乘法按“行乘列”定义：
\[
(\ten A\ten B)_{ij}=\sum_{k=1}^{3}A_{ik}B_{kj}=A_{ik}B_{kj}.
\]
矩阵乘法满足结合律：
\[
(\ten A\ten B)\ten C=\ten A(\ten B\ten C).
\]
本片段在 §1.5 把因子从左边乘进去、重排括号时直接用这条规则；P7 详解在用到的那一步再说明它的分量依据。单位矩阵 $\unitd$ 满足 $\unitd\vect x=\vect x$，其分量为 $\delta_{ij}$。

\textbf{转置。}转置矩阵 $\ten T^{\mathsf T}$ 的分量是
\[
(T^{\mathsf T})_{ij}=T_{ji}.
\]
转置的关键作用是把它从点积的一边搬到另一边：
\[
(\ten T\vect x)\cdot\vect y=\vect x\cdot(\ten T^{\mathsf T}\vect y).
\]
证明只是一次指标改名：
\[
(\ten T\vect x)\cdot\vect y
=\sum_{i=1}^{3}\left(\sum_{j=1}^{3}T_{ij}x_j\right)y_i
=\sum_{i,j}T_{ij}x_jy_i,
\]
而
\[
\vect x\cdot(\ten T^{\mathsf T}\vect y)
=\sum_{j=1}^{3}x_j\left(\sum_{i=1}^{3}T_{ij}y_i\right)
=\sum_{i,j}T_{ij}x_jy_i .
\]
两式相同。这个恒等式在 §1.4 证明“不同特征值的特征向量正交”时是关键工具。

\textbf{对称。}若 $\ten T^{\mathsf T}=\ten T$，即 $T_{ij}=T_{ji}$，就说 $\ten T$ 是\textbf{实对称矩阵}。本节的主角就是它。注意“对称”是矩阵分量的性质，不是“矩阵看起来左右对称”的视觉判断；换一组基后，同一个线性映射的矩阵可能不再对称。本节始终在正交归一坐标里讨论；若新基仍是正交归一基，对称性会保持，这一点等 §1.5 定义了换基矩阵 $\ten R$ 之后再算。

\textbf{小例。}取
\[
\ten A=\begin{pmatrix}1&2\\3&4\end{pmatrix},\qquad
\vect x=\begin{pmatrix}1\\0\end{pmatrix},\qquad
\vect y=\begin{pmatrix}0\\1\end{pmatrix}.
\]
则 $\ten A\vect x=(1,3)$，所以 $(\ten A\vect x)\cdot\vect y=3$。另一方面
\[
\ten A^{\mathsf T}\vect y=\begin{pmatrix}1&3\\2&4\end{pmatrix}\begin{pmatrix}0\\1\end{pmatrix}
=\begin{pmatrix}3\\4\end{pmatrix},
\qquad
\vect x\cdot(\ten A^{\mathsf T}\vect y)=3 .
\]
两边相同。这个例子特意用了不对称的 $\ten A$，因为恒等式对任何矩阵都成立，不需要对称。

\subhdr{0.3\quad 行列式与非零解判据}

$2\times2$ 行列式是
\[
\det\begin{pmatrix}a&b\\c&d\end{pmatrix}=ad-bc .
\]
$3\times3$ 行列式按第一行展开：
\[
\det\begin{pmatrix}
a_{11}&a_{12}&a_{13}\\
a_{21}&a_{22}&a_{23}\\
a_{31}&a_{32}&a_{33}
\end{pmatrix}
=a_{11}(a_{22}a_{33}-a_{23}a_{32})
-a_{12}(a_{21}a_{33}-a_{23}a_{31})
+a_{13}(a_{21}a_{32}-a_{22}a_{31}).
\]
本节只需要它的一条基本性质：

\begin{center}
\fcolorbox{violet}{bandviolet}{\parbox{0.94\linewidth}{\vspace{3pt}
\textbf{非零解判据。}设 $\ten A$ 是 $3\times3$ 实矩阵。齐次方程 $\ten A\vect x=\vect 0$ 有非零解，当且仅当
\[
\det\ten A=0 .
\]
\vspace{3pt}}}
\end{center}

这条判据的“$\det\ne0$”一半可以用逆矩阵说明：若 $\det\ten A\ne0$，则 $\ten A$ 可逆，$\ten A\vect x=\vect 0$ 两边左乘 $\ten A^{-1}$ 得 $\vect x=\vect 0$。反过来，若 $\det\ten A=0$，用消元把 $\ten A$ 化成上三角；三个对角元不可能全非零（否则行列式非零），于是回代时出现自由未知量，可以取非零解。这里用的是简单线性代数里的消元事实，不引入秩或维数的新语言。这条判据在 §1.3 建立特征方程、§1.4 从特征值求特征向量时都会用到。

\textbf{迹。}矩阵对角线元素之和叫\textbf{迹}：
\[
\tr\ten T=T_{11}+T_{22}+T_{33}.
\]
§1.3 用它抽查特征值之和；P7 会进一步证明它在正交换基下不变。

\subhdr{0.4\quad 奇次实多项式一定有实根}

§1.4 需要一个存在性事实：一个三次实系数多项式至少有一个实根。下面的引理把这件事说清楚；§1.3 定义的特征多项式正是这种多项式。

\begin{center}
\fcolorbox{violet}{bandviolet}{\parbox{0.94\linewidth}{\vspace{3pt}
\textbf{引理（奇次实多项式有实根）。}设
\[
p(t)=a_nt^n+a_{n-1}t^{n-1}+\cdots+a_0,\qquad a_n\ne0,
\]
其中 $n$ 是奇数，系数都是实数。则存在实数 $\lambda$ 使 $p(\lambda)=0$。
\vspace{3pt}}}
\end{center}

\textbf{证明。}若 $a_n<0$，把 $p$ 换成 $-p$；根不变，所以不妨设 $a_n>0$。当 $t>0$ 时，
\[
p(t)=t^n\left(a_n+\frac{a_{n-1}}{t}+\cdots+\frac{a_0}{t^n}\right).
\]
当 $t$ 足够大，括号内的量与 $a_n$ 的距离小于 $a_n/2$，括号为正，故 $p(t)>0$。当 $t<0$ 且 $|t|$ 足够大，$t^n<0$（因为 $n$ 为奇数），括号仍趋近于 $a_n>0$，故 $p(t)<0$。于是可以取到实数 $A<B$，使
\[
p(A)<0<p(B).
\]

多项式 $p$ 是连续函数。连续函数在闭区间 $[A,B]$ 上有介值性质：若两端函数值异号，则中间某点函数值为零。为看清这条性质，用二分法：令 $A_0=A$，$B_0=B$。取中点 $m_k$；若 $p(m_k)=0$，已经找到根；否则保留使 $p$ 在两端异号的那一半，记为 $[A_{k+1},B_{k+1}]$。这样得到嵌套区间
\[
[A_0,B_0]\supset[A_1,B_1]\supset[A_2,B_2]\supset\cdots,
\qquad
B_k-A_k=\frac{B-A}{2^k}\to0 .
\]
实数的完备性保证这些区间有唯一公共点 $\lambda$。因为 $p$ 连续，$p(A_k)$ 与 $p(B_k)$ 都趋于 $p(\lambda)$；而它们始终异号，所以 $p(\lambda)$ 既不能为正也不能为负，只能为 $0$。证完。

\begin{center}
\fcolorbox{rulepurple}{bandcream}{\parbox{0.94\linewidth}{\vspace{3pt}
\textbf{这条引理服务哪一步。}§1.4 第二步要证明：$3\times3$ 实对称矩阵至少有一个实特征值。§1.3 定义的特征多项式是三次实系数多项式，由本引理得到实根 $\lambda_1$；再由 §0.3 的非零解判据得到特征向量。证明谱定理的存在性时，这是唯一的“分析”输入；后面的 $2\times2$ 情形完全是代数计算。
\vspace{3pt}}}
\end{center}

\hdr{§1}{二次型与主轴：椭球、特征值、谱定理}

现在回到本节的主问题。给定一个 $3\times3$ 实矩阵 $\ten T$，二次型
\[
E(\vect a)=\vect a\cdot\ten T\vect a
\]
在每一个方向 $\vect a$ 上给出一个数。我们想找一组互相垂直的坐标轴，使 $E$ 在这组轴上没有交叉项；如果找得到，这些轴就是\textbf{主轴}，轴上的系数就是\textbf{特征值}。本节的主线是：

\begin{enumerate}[label=\arabic*.,leftmargin=1.8em,itemsep=2pt]
\item 二次型只看见 $\ten T$ 的对称部分，所以可以假设 $\ten T$ 对称；
\item 在二维例子里，主轴让椭圆变正，交叉项消失，系数正好是特征值；
\item 特征值与特征向量怎么算；
\item 对任何实对称 $\ten T$，这种主轴一定存在——这就是谱定理，本节给出完整证明；
\item 用谱定理把二次型写成三个平方项，并说明惯量、应力怎样接上。
\end{enumerate}

\subhdr{1.1\quad 二次型与它的对称部分（必读）}

把 $\ten T$ 的分量记作 $T_{ik}$，$\vect a=(a_1,a_2,a_3)$。二次型写成指标和就是
\[
E(\vect a)=\vect a\cdot\ten T\vect a
=\sum_{i=1}^{3}\sum_{k=1}^{3}T_{ik}a_ia_k .
\]
写开就是
\[
E=T_{11}a_1^2+T_{22}a_2^2+T_{33}a_3^2
+(T_{12}+T_{21})a_1a_2
+(T_{13}+T_{31})a_1a_3
+(T_{23}+T_{32})a_2a_3 .
\]
请注意交叉项的系数是 $T_{12}+T_{21}$，不是无条件的一个 $2T_{12}$。只有当 $\ten T$ 对称，即 $T_{12}=T_{21}$ 时，才可以把它写成 $2T_{12}a_1a_2$。这个差别看起来只是记账，实际上决定了“二次型只与对称部分有关”这条结论。

\textbf{为什么只与对称部分有关。}在求和
\[
E=\sum_{i,k}T_{ik}a_ia_k
\]
里，$i,k$ 都只是求和的哑指标。把两者的名字对换（只是换字母，值不变）：
\[
\sum_{i,k}T_{ik}a_ia_k
=\sum_{i,k}T_{ki}a_ka_i
=\sum_{i,k}T_{ki}a_ia_k .
\]
把原来的式子与这个式子相加再除以 $2$，得到
\[
E=\sum_{i,k}\frac{T_{ik}+T_{ki}}{2}\,a_ia_k
=\sum_{i,k}S_{ik}a_ia_k,
\qquad
S_{ik}=\frac{T_{ik}+T_{ki}}{2}.
\]
矩阵
\[
\ten S=\frac12\bigl(\ten T+\ten T^{\mathsf T}\bigr)
\]
叫 $\ten T$ 的\textbf{对称部分}。上面的计算说明：把 $\ten T$ 换成 $\ten S$，二次型 $E$ 一字不变。所以从这以后，凡讲二次型，都默认矩阵已经对称；如果题目给了一个不对称的 $\ten T$，先取 $\ten S$ 再谈主轴。

\textbf{教学小例。}取
\[
\ten T=\begin{pmatrix}1&2\\0&1\end{pmatrix},
\qquad
\vect a=\begin{pmatrix}x\\y\end{pmatrix}.
\]
直接算：
\[
E=\begin{pmatrix}x&y\end{pmatrix}
\begin{pmatrix}1&2\\0&1\end{pmatrix}
\begin{pmatrix}x\\y\end{pmatrix}
=x^2+2xy+y^2=(x+y)^2 .
\]
对称部分是
\[
\ten S=\frac12\left[
\begin{pmatrix}1&2\\0&1\end{pmatrix}
+\begin{pmatrix}1&0\\2&1\end{pmatrix}\right]
=\begin{pmatrix}1&1\\1&1\end{pmatrix}.
\]
用 $\ten S$ 算：
\[
E=\begin{pmatrix}x&y\end{pmatrix}
\begin{pmatrix}1&1\\1&1\end{pmatrix}
\begin{pmatrix}x\\y\end{pmatrix}
=x^2+xy+yx+y^2=x^2+2xy+y^2 .
\]
结果相同。这个例子特意保留了一个不对称的 $\ten T$，因为要看清：不是 $\ten T$ 的每个分量都不起作用，而是它的\textbf{反对称部分}在二次型里被抵消了。

\begin{center}
\fcolorbox{rulepurple}{bandcream}{\parbox{0.94\linewidth}{\vspace{3pt}
\textbf{不要把两件事混在一起。}“二次型只看见对称部分”说的是 $E(\vect a)=\vect a\cdot\ten T\vect a$ 这个数。$\ten T$ 作为线性算子作用在矢量上时，反对称部分一般是有作用的；例如叉乘就可以写成一个反对称矩阵的作用。本片段只讨论二次型，所以后面总把矩阵取成对称。
\vspace{3pt}}}
\end{center}

\textbf{物理上它长什么样。}刚体的转动动能是
\[
K=\frac12\vect\omega\cdot\ten J\vect\omega
=\frac12J_{ik}\omega_i\omega_k,
\]
正是一个二次型（差一个因子 $\frac12$）。惯量矩阵 $\ten J$ 本来对称，所以“只取对称部分”在这里不损失任何信息。§1.6 会把这个对应说完整；现在只需要记住：能量是 $\vect\omega$ 的二次函数，而二次函数的系数天然是一张 $3\times3$ 的表。

\subhdr{1.2\quad 等值面、椭球与主轴（必读）}

先把主轴说准。\textbf{定义（特征值与特征向量）。}若存在非零矢量 $\vect n$ 和数 $\lambda$ 使
\[
\ten T\vect n=\lambda\vect n,
\]
就称 $\lambda$ 是 $\ten T$ 的\textbf{特征值}，$\vect n$ 是对应的\textbf{特征向量}；$\vect n$ 张成的直线叫一条\textbf{主轴方向}。

几何上，$\ten T$ 把 $\vect n$ 张成的直线仍送到这条直线上：沿这条直线取带符号坐标 $t$（点是 $t\vect n$），则
\[
\ten T(t\vect n)=t\,\ten T\vect n=(\lambda t)\vect n,
\]
所以带符号坐标被乘以 $\lambda$，长度被乘以 $|\lambda|$。若 $\lambda>0$，方向不变；若 $\lambda<0$，方向反向；若 $\lambda=0$，整条直线被压到原点，$\vect n$ 仍按定义是特征向量，但 $\ten T\vect n$ 没有方向。因此准确地说：$\ten T$ 保持这条直线不变，沿它的带符号坐标乘以 $\lambda$；只有 $\lambda\ge0$ 时才能把“长度乘以 $\lambda$”当字面话，一般情况下长度乘以的是 $|\lambda|$。

现在令 $E(\vect a)=1$，得到一个曲面。先用二维看没有交叉项的情形：曲面是正放的椭圆；三维的椭球完全同理。若
\[
\ten T=\begin{pmatrix}1&0\\0&4\end{pmatrix},
\]
则
\[
E=x^2+4y^2=1
\quad\Longleftrightarrow\quad
\frac{x^2}{1^2}+\frac{y^2}{(1/2)^2}=1 .
\]
这是半轴分别为 $1$ 与 $1/2$ 的椭圆。两个对角元 $1,4$ 正好是 $\ten T$ 的两个特征值（直接看 $\ten T\vect e_1=\vect e_1$、$\ten T\vect e_2=4\vect e_2$）；沿第一个轴的半轴是 $1/\sqrt1=1$，沿第二个轴的半轴是 $1/\sqrt4=1/2$。一句话：\textbf{在 $\lambda>0$ 时，半轴等于 $1/\sqrt{\lambda}$，特征值越大，那个方向越短。}

\textbf{有交叉项的例子。}取
\[
\ten T=\begin{pmatrix}2&1\\1&2\end{pmatrix},
\qquad
E=2x^2+2xy+2y^2=1 .
\]
这个椭圆是斜的。先直接验证两个方向：
\[
\ten T\begin{pmatrix}1\\1\end{pmatrix}
=\begin{pmatrix}3\\3\end{pmatrix}
=3\begin{pmatrix}1\\1\end{pmatrix},
\qquad
\ten T\begin{pmatrix}-1\\1\end{pmatrix}
=\begin{pmatrix}-1\\1\end{pmatrix}
=1\begin{pmatrix}-1\\1\end{pmatrix}.
\]
所以 $\vect n_1=(1,1)/\sqrt2$ 对应特征值 $3$，$\vect n_2=(-1,1)/\sqrt2$ 对应特征值 $1$，两个方向正交。令
\[
u=\vect n_1\cdot(x,y)=\frac{x+y}{\sqrt2},
\qquad
v=\vect n_2\cdot(x,y)=\frac{-x+y}{\sqrt2}.
\]
反解是
\[
x=\frac{u-v}{\sqrt2},\qquad y=\frac{u+v}{\sqrt2}.
\]
代入：
\[
x^2+y^2=u^2+v^2,\qquad
xy=\frac{u^2-v^2}{2},
\]
所以
\[
E=2(x^2+y^2)+2xy
=2(u^2+v^2)+(u^2-v^2)
=3u^2+v^2 .
\]
交叉项真的消失了。在新的 $u,v$ 轴上，二次型是 $3u^2+v^2$；沿 $u$ 的半轴是 $1/\sqrt3$，沿 $v$ 的半轴是 $1$。让交叉项消失的那组坐标轴，就是\textbf{主轴}。

\needspace{0.34\textheight}
\begin{center}
\begin{minipage}{0.94\linewidth}
\centering
\begin{tikzpicture}[scale=2.35,>={Stealth[length=1.6mm]}]
  \begin{scope}[rotate=45]
    \fill[bandviolet,opacity=0.55] (0,0) ellipse ({0.5774} and {1});
    \draw[violet,thick] (0,0) ellipse ({0.5774} and {1});
  \end{scope}
  \draw[inkgray,densely dotted,->] (-1.25,0)--(1.30,0) node[right]{\scriptsize $x$};
  \draw[inkgray,densely dotted,->] (0,-1.20)--(0,1.30) node[above]{\scriptsize $y$};
  \draw[violet,thick,->] (0,0)--(0.72,0.72) node[above right]{\scriptsize $u$（$\vect n_1$）};
  \draw[violet,thick,->] (0,0)--(-0.86,0.86) node[above left]{\scriptsize $v$（$\vect n_2$）};
  \node[violet] at (0.56,0.29) {\scriptsize $\lambda_1=3$};
  \node[violet] at (-0.57,0.29) {\scriptsize $\lambda_2=1$};
  \fill (0,0) circle (0.7pt);
\end{tikzpicture}\\[2pt]
\figcap{\mbox{\textbf{图 1}（自绘）：}$E=2x^2+2xy+2y^2=1$ 的等值线。灰虚线是原来的 $x,y$ 轴；紫线是主轴 $u=\vect n_1$、$v=\vect n_2$。沿 $u$ 的特征值是 $3$，半轴只有 $1/\sqrt3\approx0.577$；沿 $v$ 的特征值是 $1$，半轴是 $1$。图中的椭圆长短轴与这两个数对应：\textbf{半轴 $=1/\sqrt{\lambda}$，特征值越大那一轴越短}。}
\end{minipage}
\end{center}

这张图与上面的代入逐项对应：$v=0$ 时 $(x,y)=u(1,1)/\sqrt2$，正是 $\vect n_1$ 方向；$u=0$ 时 $(x,y)=v(-1,1)/\sqrt2$，正是 $\vect n_2$ 方向。图中的 $u$ 轴沿 $45^\circ$，$v$ 轴沿 $135^\circ$，椭圆在 $v$ 方向更长，因为 $v$ 方向的半轴是 $1$，而 $u$ 方向的半轴只有 $0.577$。

\begin{center}
\fcolorbox{rulepurple}{bandcream}{\parbox{0.94\linewidth}{\vspace{3pt}
\textbf{边界情形。}上面的“半轴 $=1/\sqrt{\lambda}$”是在 $\lambda>0$ 时说的。若某个特征值为负，$E=1$ 在该方向上没有实点，曲面可能是双曲线或空集；若特征值为零，曲面退化（例如一对平行线或柱面）。但“主轴方向是特征向量、主值是特征值”仍然成立。谱定理只要求 $\ten T$ 实对称，不要求它正定；正定只保证 $E=1$ 是真椭球。惯量矩阵通常正定或半正定，应力矩阵则可能有负特征值——那时仍用主轴与主应力，但不画椭球。
\vspace{3pt}}}
\end{center}

\subhdr{1.3\quad 特征值与特征向量怎么算（必读）}

把 $\ten T\vect n=\lambda\vect n$ 移项：
\[
(\ten T-\lambda\unitd)\vect n=\vect 0 .
\]
要它有非零解，由 §0.3 的非零解判据，必须
\[
\det(\ten T-\lambda\unitd)=0 .
\]
左边是 $\lambda$ 的多项式，叫\textbf{特征多项式}；这个方程叫\textbf{特征方程}。对 $3\times3$ 矩阵，按第一行展开后它是三次多项式；对 $2\times2$ 矩阵，直接算得
\[
\det\begin{pmatrix}a-\lambda&b\\c&d-\lambda\end{pmatrix}
=\lambda^2-(a+d)\lambda+(ad-bc),
\]
这个二次公式在 §1.4 的 $2\times2$ 引理里会直接用。解出特征值 $\lambda$ 后，再把它代回 $(\ten T-\lambda\unitd)\vect n=\vect 0$，解出的非零 $\vect n$ 就是特征向量。

\textbf{三维教学示范。}取
\[
\ten D=\begin{pmatrix}6&2&0\\2&6&0\\0&0&1\end{pmatrix}.
\]
先算特征多项式：
\[
\det(\ten D-\lambda\unitd)
=(1-\lambda)\bigl[(6-\lambda)^2-4\bigr]
=(1-\lambda)(\lambda-4)(\lambda-8).
\]
所以特征值是
\[
\lambda_1=8,\qquad \lambda_2=4,\qquad \lambda_3=1.
\]
再求方向：
\begin{itemize}[leftmargin=1.8em,itemsep=2pt]
\item $\lambda=8$：$(\ten D-8\unitd)\vect n=\vect0$ 的前两行给出 $-2x+2y=0$，即 $x=y$；第三行给出 $-7z=0$，即 $z=0$。取
\[
\vect n_1=\frac{1}{\sqrt2}\begin{pmatrix}1\\1\\0\end{pmatrix}.
\]
\item $\lambda=4$：前两行给出 $2x+2y=0$，即 $y=-x$；第三行给出 $3z=0$。取
\[
\vect n_2=\frac{1}{\sqrt2}\begin{pmatrix}1\\-1\\0\end{pmatrix}.
\]
\item $\lambda=1$：前两行是 $5x+2y=0$ 与 $2x+5y=0$，两式合起来只能 $x=y=0$；$z$ 任意。取
\[
\vect n_3=\begin{pmatrix}0\\0\\1\end{pmatrix}.
\]
\end{itemize}
三个方向两两正交，长度都是 $1$。把它们按“列乘行”的方式拼回去：
\[
8\,\vect n_1\vect n_1^{\mathsf T}
+4\,\vect n_2\vect n_2^{\mathsf T}
+1\,\vect n_3\vect n_3^{\mathsf T}
\]
\[
=8\cdot\frac12\begin{pmatrix}1&1&0\\1&1&0\\0&0&0\end{pmatrix}
+4\cdot\frac12\begin{pmatrix}1&-1&0\\-1&1&0\\0&0&0\end{pmatrix}
+\begin{pmatrix}0&0&0\\0&0&0\\0&0&1\end{pmatrix}
=\begin{pmatrix}6&2&0\\2&6&0\\0&0&1\end{pmatrix}.
\]
最后一步逐项相加即可。这里 $\vect n\vect n^{\mathsf T}$ 是“列乘行”得到的 $3\times3$ 矩阵，$(i,j)$ 元是 $n_in_j$；它作用在 $\vect x$ 上给出 $\vect n(\vect n\cdot\vect x)$，也就是把 $\vect x$ 投影到 $\vect n$ 方向。§1.5 会把这个拼法写成一般公式。

\textbf{重复特征值。}例如
\[
\ten D=\operatorname{diag}(2,2,5).
\]
$\lambda=2$ 对应整个 $xy$ 平面：平面里任意非零矢量都满足 $\ten D\vect n=2\vect n$，所以任意一对正交单位矢量都可以当主轴；$\lambda=5$ 对应 $z$ 轴。重复特征值不破坏谱定理，只让“主轴的选择”不唯一。这一点在惯量张量的圆柱对称、对称陀螺里都会出现。

\textbf{用迹和行列式抽查。}对上面的 $\ten D$，
\[
\tr\ten D=6+6+1=13=8+4+1,\qquad
\det\ten D=32=8\cdot4\cdot1 .
\]
这两个等式对任何矩阵都成立（特征值之和等于迹，特征值之积等于行列式；比较特征多项式按第一行展开的系数即可看出），可以用来抓符号错误。但它们只是\textbf{检查}，不是谱定理的证明；证明在 §1.4。

\subhdr{1.4\quad 谱定理：陈述与完整证明（必读）}

\begin{center}
\fcolorbox{violet}{bandviolet}{\parbox{0.94\linewidth}{\vspace{4pt}
\textbf{谱定理（实对称矩阵，三维）。}设 $\ten T$ 是 $3\times3$ 实矩阵，且 $\ten T^{\mathsf T}=\ten T$。则：
\begin{enumerate}[label=(\roman*),leftmargin=2.0em,itemsep=3pt]
\item $\ten T$ 的三个特征值（按重数计）都是实数；
\item 属于不同特征值的特征向量互相正交；
\item 存在三个两两正交的单位矢量 $\vect n_1,\vect n_2,\vect n_3$ 和实数 $\lambda_1,\lambda_2,\lambda_3$，使
\[
\ten T\vect n_k=\lambda_k\vect n_k\quad(k=1,2,3),
\]
并且
\[
\ten T=\lambda_1\vect n_1\vect n_1^{\mathsf T}
+\lambda_2\vect n_2\vect n_2^{\mathsf T}
+\lambda_3\vect n_3\vect n_3^{\mathsf T}.
\]
等价地说：以 $\{\vect n_1,\vect n_2,\vect n_3\}$ 为基，$\ten T$ 的矩阵是对角阵 $\operatorname{diag}(\lambda_1,\lambda_2,\lambda_3)$。
\end{enumerate}
\vspace{4pt}}}
\end{center}

这里的 $\vect n_k\vect n_k^{\mathsf T}$ 就是 §1.3 用过的“列乘行”矩阵；它不是装饰写法，而是一个真正的 $3\times3$ 矩阵。定理说：任何实对称矩阵都可以写成三个这样的投影矩阵的加权和，权重就是特征值。

\textbf{证明路线。}证明分五步：
\begin{enumerate}[label=\arabic*.,leftmargin=1.8em,itemsep=2pt]
\item 先完全解决 $2\times2$ 情形；
\item 用 §0.4 证明 $3\times3$ 实对称矩阵至少有一个实特征值；
\item 证明 $\ten T$ 保持这个特征向量的正交补，把问题降成 $2\times2$；
\item 用第 1 步的 $2\times2$ 引理得到另外两个正交特征向量；
\item 验证分解式，并补证“不同特征值的特征向量正交”。
\end{enumerate}
这样安排的原因很实际：$3\times3$ 的直接计算会遇到三次方程，而三次方程没有简单的求根公式；但一旦有了一个特征向量，剩下的平面就是二维的，$2\times2$ 情形可以用判别式完全算清楚。

\step{第一步：$2\times2$ 引理。}
\begin{center}
\fcolorbox{violet}{bandviolet}{\parbox{0.94\linewidth}{\vspace{3pt}
\textbf{引理（$2\times2$ 实对称矩阵）。}设
\[
\ten B=\begin{pmatrix}a&b\\b&d\end{pmatrix},\qquad a,b,d\in\R .
\]
则 $\ten B$ 有两个实特征值（可能相等），并且存在一组正交归一基，使 $\ten B$ 在这组基下的矩阵是对角阵。
\vspace{3pt}}}
\end{center}

\textbf{证明。}特征多项式是
\[
p_{\ten B}(\lambda)=\det(\ten B-\lambda\unitd)
=\lambda^2-(a+d)\lambda+(ad-b^2).
\]
判别式
\[
\Delta=(a+d)^2-4(ad-b^2)
=(a-d)^2+4b^2\ge0 .
\]
所以两个根
\[
\mu_\pm=\frac{(a+d)\pm\sqrt{\Delta}}{2}
\]
都是实数。

先处理 $b=0$ 的情形。此时 $\ten B=\operatorname{diag}(a,d)$，特征值是 $a,d$，对应的特征向量可以取 $\vect e_1,\vect e_2$。若 $a=d$，则 $\ten B=a\unitd$，任意单位矢量都是特征向量，任取一组正交归一基即可。

再处理 $b\ne0$ 的情形。此时 $\Delta=(a-d)^2+4b^2>0$，两个根 $\mu_+\ne\mu_-$。对每个根 $\mu$，令
\[
\vect v_\mu=\begin{pmatrix}b\\ \mu-a\end{pmatrix}.
\]
因为 $b\ne0$，$\vect v_\mu\ne\vect0$。验证它确实是特征向量：
\[
\ten B\vect v_\mu
=\begin{pmatrix}a&b\\b&d\end{pmatrix}
\begin{pmatrix}b\\ \mu-a\end{pmatrix}
=\begin{pmatrix}ab+b(\mu-a)\\ b^2+d(\mu-a)\end{pmatrix}
=\begin{pmatrix}b\mu\\ b^2+d\mu-ad\end{pmatrix}.
\]
由 $p_{\ten B}(\mu)=0$ 得
\[
\mu^2-(a+d)\mu+(ad-b^2)=0
\quad\Longrightarrow\quad
\mu^2-a\mu=d\mu-ad+b^2 .
\]
而第二个分量正好是
\[
b^2+d\mu-ad=\mu^2-a\mu=\mu(\mu-a),
\]
所以 $\ten B\vect v_\mu=\mu\vect v_\mu$。

最后看两个特征向量是否正交。利用根与系数的关系
\[
\mu_++\mu_-=a+d,\qquad
\mu_+\mu_-=ad-b^2,
\]
得到
\[
(\mu_+-a)(\mu_--a)
=\mu_+\mu_- -a(\mu_++\mu_-)+a^2
=(ad-b^2)-a(a+d)+a^2
=-b^2 .
\]
于是
\[
\vect v_+ \cdot \vect v_-
=b^2+(\mu_+-a)(\mu_--a)
=b^2-b^2=0 .
\]
把 $\vect v_\pm$ 分别单位化，就得到一组正交归一特征基。$2\times2$ 引理证完。

\step{第二步：$3\times3$ 实对称矩阵至少有一个实特征值。}
设 $\ten T$ 是 $3\times3$ 实对称矩阵。它的特征多项式
\[
p_{\ten T}(\lambda)=\det(\ten T-\lambda\unitd)
\]
按第一行展开后是三次实系数多项式，最高次项为 $-\lambda^3$。由 §0.4 的引理，存在实数 $\lambda_1$ 使 $p_{\ten T}(\lambda_1)=0$。再由 §0.3 的非零解判据，方程
\[
(\ten T-\lambda_1\unitd)\vect n=\vect0
\]
有非零解；取一个这样的解，单位化后记作 $\vect u_1$。于是
\[
\ten T\vect u_1=\lambda_1\vect u_1,\qquad |\vect u_1|=1 .
\]

\step{第三步：$\ten T$ 保持 $\vect u_1$ 的正交补。}
令
\[
P=\{\vect x\in\R^3:\vect x\cdot\vect u_1=0\}.
\]
这是过原点、与 $\vect u_1$ 垂直的平面。若 $\vect x\in P$，则用 §0.2 的转置恒等式和 $\ten T^{\mathsf T}=\ten T$：
\[
(\ten T\vect x)\cdot\vect u_1
=\vect x\cdot(\ten T^{\mathsf T}\vect u_1)
=\vect x\cdot(\ten T\vect u_1)
=\lambda_1(\vect x\cdot\vect u_1)=0 .
\]
所以 $\ten T\vect x$ 仍在 $P$ 里，即 $\ten T(P)\subset P$。这一步是降维的关键：$\ten T$ 不会把平面 $P$ 里的矢量“翘出”平面。

\begin{center}
\begin{tikzpicture}[scale=1.02,>={Stealth[length=1.6mm]}]
  \fill[bandviolet,opacity=0.55] (-1.65,-0.45)--(1.20,-0.45)--(1.75,0.55)--(-1.15,0.55)--cycle;
  \draw[rulepurple] (-1.65,-0.45)--(1.20,-0.45)--(1.75,0.55)--(-1.15,0.55)--cycle;
  \draw[thick,->] (0,0)--(0,2.0) node[above]{\scriptsize $\vect u_1$};
  \draw[violet,thick,->] (0,0)--(0.95,-0.02) node[below]{\scriptsize $\vect x$};
  \draw[violet,thick,->] (0,0)--(0.30,0.38) node[above right]{\scriptsize $\ten T\vect x$};
  \node at (1.25,0.28) {\scriptsize $P=\vect u_1^{\perp}$};
  \fill (0,0) circle (0.6pt);
\end{tikzpicture}\\[2pt]
\figcap{\mbox{\textbf{图 2}（自绘，示意）：}$\vect u_1$ 与平面 $P$ 垂直；若 $\vect x\in P$，则 $\ten T\vect x$ 仍在 $P$ 内。图中只画方向与包含关系，不表示长度；它对应第三步的等式 $(\ten T\vect x)\cdot\vect u_1=0$。}
\end{center}

在 $P$ 内取一组正交归一基 $\vect e_2,\vect e_3$。这样的基一定存在：先取 $P$ 内任意非零矢量并单位化，记为 $\vect e_2$；再取 $P$ 内一个与 $\vect e_2$ 不平行的矢量 $\vect w$（平面是二维的，这样的 $\vect w$ 存在），令
\[
\vect z=\vect w-(\vect w\cdot\vect e_2)\vect e_2 .
\]
由 §0.1 的投影公式，$\vect z\cdot\vect e_2=0$，且 $\vect z\ne\vect0$（否则 $\vect w$ 与 $\vect e_2$ 平行，矛盾）。把 $\vect z$ 单位化就得到 $\vect e_3$。因为 $\ten T(P)\subset P$，每个 $\ten T\vect e_j$ 都能写成 $\vect e_2,\vect e_3$ 的线性组合。定义
\[
B_{ij}=\vect e_i\cdot\ten T\vect e_j,\qquad i,j=2,3,
\]
则
\[
\ten T\vect e_j=\sum_{i=2}^{3}B_{ij}\vect e_i .
\]
矩阵 $\ten B=(B_{ij})$ 是 $2\times2$ 实矩阵，而且对称：
\[
B_{ij}=\vect e_i\cdot\ten T\vect e_j
=(\ten T\vect e_i)\cdot\vect e_j
=\vect e_j\cdot\ten T\vect e_i
=B_{ji},
\]
其中第二个等号用了 $\ten T$ 的对称性，第三个等号用了点积的交换律。于是 $\ten B$ 是 $2\times2$ 实对称矩阵。

\step{第四步：在平面内用 $2\times2$ 引理。}

先用一个具体对象看清这一步在做什么。取标准基 $\vect i=(1,0,0)$、$\vect j=(0,1,0)$、$\vect k=(0,0,1)$，并令
\[
\ten T=\begin{pmatrix}5&0&0\\0&2&1\\0&1&2\end{pmatrix}.
\]
它实对称，$\vect u_1=\vect i$ 是特征向量，$\ten T\vect u_1=5\vect u_1$，所以 $\lambda_1=5$，而 $P=\{\vect x:\vect x\cdot\vect u_1=0\}$ 就是 $yz$ 平面。在 $P$ 内取 $\vect e_2=\vect j$、$\vect e_3=\vect k$，则第三步的 $\ten B$ 是
\[
B_{ij}=\vect e_i\cdot\ten T\vect e_j
\quad\Longrightarrow\quad
\ten B=\begin{pmatrix}2&1\\1&2\end{pmatrix}.
\]
这个 $\ten B$ 正是 §1.2 用过的 $2\times2$ 矩阵。$P$ 中的三维矢量 $y\vect e_2+z\vect e_3$ 用两个坐标 $(y,z)$ 表示，$\ten T$ 在 $P$ 上的作用由 $\ten B$ 完成：
\[
\ten T(y\vect e_2+z\vect e_3)
=(2y+z)\vect e_2+(y+2z)\vect e_3,
\]
右端的两个系数正是 $\ten B(y,z)^{\mathsf T}$。例如 $\ten B(1,0)^{\mathsf T}=(2,1)^{\mathsf T}$，拼回三维就是 $\ten T\vect e_2=2\vect e_2+\vect e_3$；$\ten B(1,1)^{\mathsf T}=3(1,1)^{\mathsf T}$ 则给出特征方向。$\ten B$ 的特征向量是
\[
\vect c_2=\frac{1}{\sqrt2}\begin{pmatrix}1\\1\end{pmatrix}\ (\mu_2=3),
\qquad
\vect c_3=\frac{1}{\sqrt2}\begin{pmatrix}1\\-1\end{pmatrix}\ (\mu_3=1).
\]
把二维特征向量拼回三维：
\[
\vect u_2=c_{22}\vect e_2+c_{32}\vect e_3
=\frac{1}{\sqrt2}(\vect e_2+\vect e_3)
=\frac{1}{\sqrt2}\begin{pmatrix}0\\1\\1\end{pmatrix},
\]
\[
\vect u_3=c_{23}\vect e_2+c_{33}\vect e_3
=\frac{1}{\sqrt2}(\vect e_2-\vect e_3)
=\frac{1}{\sqrt2}\begin{pmatrix}0\\1\\-1\end{pmatrix}.
\]
直接核对：
\[
\ten T\vect u_2=3\vect u_2,\qquad
\ten T\vect u_3=\vect u_3 .
\]
所以 $\vect u_1,\vect u_2,\vect u_3$ 是一组正交归一特征基，对应的特征值是 $5,3,1$。这个例子对应一般证明的第四步：$\ten B$ 的两个特征向量 $\vect c_2,\vect c_3$ 给出两个坐标方向，按 $\vect u_k=c_{2k}\vect e_2+c_{3k}\vect e_3$ 拼回 $P$ 中的三维矢量，就得到 $\ten T$ 的另外两个特征向量。它只演示对应关系，不替代下面的字母论证。

现在回到一般情形。对 $\ten B$ 应用第一步的引理：存在 $2\times2$ 的正交归一特征基 $\vect c_2,\vect c_3$ 和实数 $\mu_2,\mu_3$，使
\[
\ten B\vect c_k=\mu_k\vect c_k,\qquad k=2,3 .
\]
把 $\vect c_k=(c_{2k},c_{3k})$ 重新拼成三维矢量
\[
\vect u_k=c_{2k}\vect e_2+c_{3k}\vect e_3,\qquad k=2,3 .
\]
因为 $\vect c_2,\vect c_3$ 正交归一，$\vect u_2,\vect u_3$ 也在 $P$ 内正交归一。并且
\[
\ten T\vect u_k
=c_{2k}\ten T\vect e_2+c_{3k}\ten T\vect e_3
=\sum_{i=2}^{3}(\ten B\vect c_k)_i\vect e_i
=\mu_k\vect u_k .
\]
所以 $\vect u_2,\vect u_3$ 是 $\ten T$ 的特征向量，特征值分别是 $\mu_2,\mu_3$。为记号统一，把 $\mu_2,\mu_3$ 重新记作 $\lambda_2,\lambda_3$。又因为 $\vect u_2,\vect u_3\in P$，它们都与 $\vect u_1$ 正交。于是
\[
\vect u_1,\vect u_2,\vect u_3
\]
是一组正交归一特征基。

这一步同时说明了“所有特征值都是实数”。设 $\vect x\ne\vect0$ 是 $\ten T$ 的任一特征向量，特征值为 $\lambda$，即 $\ten T\vect x=\lambda\vect x$。把 $\vect x$ 沿 $\vect u_1$ 与平面 $P$ 分解：
\[
\vect x=\alpha\vect u_1+\vect p,\qquad
\alpha=\vect x\cdot\vect u_1,\qquad
\vect p=\vect x-\alpha\vect u_1\in P .
\]
因为 $\ten T\vect u_1=\lambda_1\vect u_1$，且第三步已证 $\ten T\vect p\in P$，
\[
\ten T\vect x=\alpha\lambda_1\vect u_1+\ten T\vect p .
\]
左边又等于 $\lambda\vect x=\lambda\alpha\vect u_1+\lambda\vect p$。比较两边的 $\vect u_1$ 分量和 $P$ 内分量，得到
\[
\alpha\lambda_1=\lambda\alpha,\qquad
\ten T\vect p=\lambda\vect p .
\]
若 $\alpha\ne0$，第一式给出 $\lambda=\lambda_1$。若 $\alpha=0$，则 $\vect x=\vect p\in P$ 且 $\vect p\ne\vect0$；把 $\vect p$ 写成 $\vect p=c_2\vect e_2+c_3\vect e_3$，第二式就是 $\ten B\vect c=\lambda\vect c$，其中 $\vect c=(c_2,c_3)\ne\vect0$，所以 $\lambda$ 是 $\ten B$ 的特征值，由第一步的 $2\times2$ 引理只能是 $\mu_2$ 或 $\mu_3$。

于是 $\ten T$ 的每个特征值都属于 $\{\lambda_1,\mu_2,\mu_3\}$；而这三个数本身都是 $\ten T$ 的特征值（分别对应 $\vect u_1,\vect u_2,\vect u_3$）。三次特征多项式按重数计有三个根，它们都是实数，所以 (i) 成立。这里不需要引入特征多项式的行列式恒等式，也不需要把 $\ten B$ 的特征多项式与原矩阵的特征多项式作比较；用到的只是“沿 $\vect u_1$ 与 $P$ 分解”这一件第三步已经建立的事。

\step{第五步：正交性与分解式。}
先补证 (ii)。设 $\ten T\vect u=\lambda\vect u$，$\ten T\vect v=\mu\vect v$，且 $\lambda\ne\mu$。则
\[
\lambda(\vect u\cdot\vect v)
=(\ten T\vect u)\cdot\vect v
=\vect u\cdot(\ten T^{\mathsf T}\vect v)
=\vect u\cdot(\ten T\vect v)
=\mu(\vect u\cdot\vect v).
\]
第二个等号是 §0.2 的转置恒等式，第三个等号用了 $\ten T^{\mathsf T}=\ten T$。移项得
\[
(\lambda-\mu)(\vect u\cdot\vect v)=0 .
\]
因为 $\lambda\ne\mu$，只能 $\vect u\cdot\vect v=0$。这就证明了 (ii)。注意这一步没有用到“特征向量已经单位化”，只用了对称性。

最后证明 (iii) 的分解式。对每个基矢量 $\vect u_j$，右边作用上去：
\[
\left(\sum_{k=1}^{3}\lambda_k\vect u_k\vect u_k^{\mathsf T}\right)\vect u_j
=\sum_{k=1}^{3}\lambda_k\vect u_k(\vect u_k\cdot\vect u_j)
=\lambda_j\vect u_j,
\]
因为 $\vect u_k\cdot\vect u_j=\delta_{kj}$。左边 $\ten T$ 作用在 $\vect u_j$ 上也给出 $\lambda_j\vect u_j$。两个线性映射在基 $\vect u_1,\vect u_2,\vect u_3$ 上作用相同，所以它们相等：
\[
\ten T=\sum_{k=1}^{3}\lambda_k\vect u_k\vect u_k^{\mathsf T}.
\]
谱定理证完。

\begin{center}
\fcolorbox{rulepurple}{bandcream}{\parbox{0.94\linewidth}{\vspace{3pt}
\textbf{这个证明实际用了什么。}它只用了：$\ten T$ 实对称；三维特征多项式是三次实系数多项式；§0.3 的非零解判据；§0.4 的实根存在性；以及 $2\times2$ 的判别式计算。它没有用正定性，没有用 $\ten T$ 可逆，也没有用多元微积分或拉格朗日乘子。因此它对惯量矩阵（通常正定或半正定）和应力矩阵（可能有负特征值）同样适用。证明路线不是唯一的；这里选“先降维到 $2\times2$”的路线，是因为它把唯一的存在性困难集中到 §0.4 一个引理上，后面的步骤都是可逐行核对的代数。
\vspace{3pt}}}
\end{center}

\textbf{重复特征值时会发生什么。}若三个特征值全相等，$\ten T=\lambda_1\unitd$：由分解式，
\[
\ten T=\lambda_1(\vect u_1\vect u_1^{\mathsf T}
+\vect u_2\vect u_2^{\mathsf T}
+\vect u_3\vect u_3^{\mathsf T})
=\lambda_1\unitd,
\]
所以任意正交归一基都是特征基。若两个特征值相等，对应的特征空间是一个平面，平面内任意一组正交归一基都可以当主轴。谱定理保证的是“存在一组”，不是“只有一组”；说“主轴方向”时，若特征值重复，必须说明是在哪个特征空间里选。

\subhdr{1.5\quad 谱定理的三个用法（必读）}

谱定理的价值不只在“存在三个主轴”，更在于它把矩阵、二次型和换基三件事同时说清楚。下面三个用法后面都会反复出现。

\textbf{用法一：换到主轴坐标系，矩阵变成对角阵。}把三个主轴按列拼成矩阵
\[
\ten R=\begin{pmatrix}\vect u_1&\vect u_2&\vect u_3\end{pmatrix}.
\]
因为三个列矢量正交归一，
\[
\ten R^{\mathsf T}\ten R=\unitd .
\]
由 $\ten T\vect u_k=\lambda_k\vect u_k$，
\[
\ten T\ten R
=\begin{pmatrix}\ten T\vect u_1&\ten T\vect u_2&\ten T\vect u_3\end{pmatrix}
=\begin{pmatrix}\lambda_1\vect u_1&\lambda_2\vect u_2&\lambda_3\vect u_3\end{pmatrix}
=\ten R\operatorname{diag}(\lambda_1,\lambda_2,\lambda_3).
\]
两边左乘 $\ten R^{\mathsf T}$，利用 $\ten R^{\mathsf T}\ten R=\unitd$，得到
\[
\ten R^{\mathsf T}\ten T\ten R
=\operatorname{diag}(\lambda_1,\lambda_2,\lambda_3).
\]
这就是“换到主轴坐标系，账本变成对角阵”的准确意思。注意 $\ten R$ 的列是新基在旧基下的分量；$\ten R^{\mathsf T}\ten T\ten R$ 正是同一线性映射在新基下的矩阵。

\textbf{正交换基保持对称。}这里 $\ten R$ 的列是一组正交归一基，所以 $\ten R^{\mathsf T}\ten R=\unitd$。若 $\ten T$ 对称，则
\[
(\ten R^{\mathsf T}\ten T\ten R)^{\mathsf T}
=\ten R^{\mathsf T}\ten T^{\mathsf T}\ten R
=\ten R^{\mathsf T}\ten T\ten R,
\]
所以 $\ten R^{\mathsf T}\ten T\ten R$ 仍是对称矩阵。这个计算放在这里，是因为直到定义了 $\ten R$ 才说得清“换基后的矩阵”是什么；它不改变 §1.4 的证明。

\textbf{用法二：二次型变成三个平方项。}任意矢量可以沿主轴展开：
\[
\vect a=\sum_{k=1}^{3}c_k\vect u_k,
\qquad
c_k=\vect u_k\cdot\vect a .
\]
因为 $\ten T\vect u_k=\lambda_k\vect u_k$，
\[
\ten T\vect a=\sum_{k=1}^{3}\lambda_k c_k\vect u_k,
\]
再与 $\vect a$ 作点积，利用 $\vect u_j\cdot\vect u_k=\delta_{jk}$：
\[
E(\vect a)=\vect a\cdot\ten T\vect a
=\sum_{k=1}^{3}\lambda_k c_k^2
=\lambda_1(\vect u_1\cdot\vect a)^2
+\lambda_2(\vect u_2\cdot\vect a)^2
+\lambda_3(\vect u_3\cdot\vect a)^2 .
\]
这就是“没有交叉项”的一般形式：交叉项消失，只剩三个投影长度的平方，各乘一个主值。

\textbf{用法三：瑞利商给出上下界。}对任意非零矢量 $\vect a$，定义
\[
\rho(\vect a)=\frac{\vect a\cdot\ten T\vect a}{\vect a\cdot\vect a}.
\]
把 $\vect a=\sum_k c_k\vect u_k$ 代入，分母是 $\sum_k c_k^2$，分子是 $\sum_k\lambda_k c_k^2$，所以
\[
\rho(\vect a)
=\frac{\sum_k\lambda_k c_k^2}{\sum_k c_k^2}
\]
是一个以 $c_k^2/\sum_j c_j^2$ 为权重的加权平均。权重非负且和为 $1$，因此
\[
\lambda_{\min}\le\rho(\vect a)\le\lambda_{\max}.
\]
其中 $\lambda_{\min}$ 与 $\lambda_{\max}$ 分别是三个特征值中最小的和最大的。这个不等式在 §1.6 说明“绕任意方向的转动惯量介于最小与最大主转动惯量之间”时直接可用。

\textbf{半轴公式。}若 $\ten T$ 正定（三个特征值都大于零），$E(\vect a)=1$ 是椭球。沿主轴 $\vect u_k$ 取 $\vect a=t\vect u_k$，其中 $t$ 是带符号的轴向坐标。则
\[
E=\lambda_k t^2=1
\quad\Longrightarrow\quad
t=\pm\frac{1}{\sqrt{\lambda_k}} .
\]
两个解是椭圆在该轴上的两个端点；半轴长是 $|t|=1/\sqrt{\lambda_k}$。若把 $t$ 就取成非负的半轴长，则 $t=1/\sqrt{\lambda_k}$，但这是先选定了“从中心沿正方向量”的那一半，不是方程自动只给出正根。若 $\lambda_k\le0$，方程没有实解，这个公式不能给出实半轴；此时曲面不是椭球，但“主轴方向仍是特征向量”不变。

\subhdr{1.6\quad 两个力学落点：惯量与应力（接口）}

这一小节只说明“谱定理怎样接上力学”，不展开惯量张量的构造、欧拉方程或应力平衡。那些内容属于后续章节；这里只把数学对应关系摆清楚。

\textbf{落点一：惯量张量。}刚体的转动惯量矩阵
\[
\ten J=\int\bigl(|\vect r|^2\unitd-\vect r\vect r^{\mathsf T}\bigr)\dd m
\]
是实对称矩阵，所以由谱定理，它有三个实特征值 $A,B,C$（叫\textbf{主转动惯量}）和三条互相垂直的主轴（叫\textbf{主转动轴}）。沿单位方向 $\vect n$ 的转动惯量是二次型
\[
I(\vect n)=\vect n\cdot\ten J\vect n .
\]
由 §1.5 的公式，
\[
I(\vect n)=A(\vect u_1\cdot\vect n)^2
+B(\vect u_2\cdot\vect n)^2
+C(\vect u_3\cdot\vect n)^2,
\]
所以 $I(\vect n)$ 介于最小与最大主转动惯量之间；当 $\vect n$ 沿某条主轴时，$I$ 正好等于对应的主转动惯量。

\textbf{小例。}取 $\ten J=\operatorname{diag}(1,2,3)$，方向 $\vect n=(1,1,1)/\sqrt3$。则
\[
I(\vect n)=\vect n\cdot\ten J\vect n
=\frac{1}{3}(1+2+3)=2,
\]
正好是三个主转动惯量的算术平均。这个例子里 $\vect n$ 不是主轴；它的 $I$ 落在最小值 $1$ 与最大值 $3$ 之间，和瑞利商结论一致。

\textbf{落点二：应力张量。}在通常的连续介质（无体力矩）里，应力张量 $\ten\sigma$ 也是实对称的，所以它有三个实主应力 $\sigma_1,\sigma_2,\sigma_3$ 和三条互相垂直的主方向。过单位法向 $\vect n$ 的截面上的面力是 $\vect t=\ten\sigma\vect n$，法向应力是 $\vect n\cdot\ten\sigma\vect n$；当 $\vect n$ 沿主方向时，面力与法向平行，这个值就是主应力。应力可以有负特征值，所以通常说“主应力”而不画椭球；数学上与惯量用的是同一个谱定理。

\begin{center}
\tabtop
\begin{tabular}{@{}p{0.19\linewidth}p{0.28\linewidth}p{0.25\linewidth}p{0.20\linewidth}@{}}
\thead{对象} & \thead{二次型} & \thead{主轴} & \thead{主值} \\
\multicolumn{4}{@{}l@{}}{\color{petal}\rule{\linewidth}{0.5pt}}\\
一般实对称 $\ten T$ & $\vect a\cdot\ten T\vect a$ & 特征向量 & 特征值 \\
惯量矩阵 $\ten J$ & $\frac12\vect\omega\cdot\ten J\vect\omega$ 或 $\vect n\cdot\ten J\vect n$ & 主转动轴 & 主转动惯量 \\
应力矩阵 $\ten\sigma$ & $\vect n\cdot\ten\sigma\vect n$ & 主方向 & 主应力 \\
\end{tabular}
\tabbot
\end{center}

这张表只是把已经讲过的对应关系并排放在一起，不是三个新定理。三行用的是同一个谱定理：实对称矩阵 $\to$ 三个实特征值 $\to$ 三条正交主轴 $\to$ 二次型没有交叉项。区别只在物理名称和“是否正定”：惯量正定时有真椭球，半正定时椭球可能退化；应力可能有负特征值，所以只谈主应力。

\textbf{回到主线。}本节的数学链条可以压缩成一句话：\textbf{二次型只看见对称部分；实对称矩阵的谱定理给出三条正交主轴；在主轴坐标系里二次型变成三个平方项，系数就是特征值。}后面的惯量、欧拉方程和应力，只要遇到“二次型 + 实对称矩阵”，就回到 §1.5 的三个公式；遇到“为什么一定有三个主轴”，就回到 §1.4 的证明。这里不再展开那些章节，是为了不让本节的主线被后续构造淹没。

\hdr{§2}{独立练习}

下面这道题只用到 §0.2 的转置与矩阵乘法、§0.3 的迹定义和 §1.5 的换基公式。请先独立完成，再翻附录；正文不提示解法。

\practice{P7}{自编概念题（承旧稿 P7）}
{设 $\ten R$ 是 $3\times3$ 正交矩阵，即 $\ten R^{\mathsf T}\ten R=\ten R\ten R^{\mathsf T}=\unitd$；$\ten T$ 是任意 $3\times3$ 实矩阵。用 $\tr(\ten A\ten B)=\tr(\ten B\ten A)$ 写出两步，证明 $\tr(\ten R^{\mathsf T}\ten T\ten R)=\tr\ten T$，并说明每一步用了什么性质。}
{正交换基下迹的不变性；循环性作为题目已给工具使用。}{ans:P7}

\clearpage
\appendix
\hdr{附录}{练习详解}

\answerstart{P7}{自编概念题（承旧稿 P7）}
{设 $\ten R$ 是 $3\times3$ 正交矩阵，即 $\ten R^{\mathsf T}\ten R=\ten R\ten R^{\mathsf T}=\unitd$；$\ten T$ 是任意 $3\times3$ 实矩阵。用 $\tr(\ten A\ten B)=\tr(\ten B\ten A)$ 写出两步，证明 $\tr(\ten R^{\mathsf T}\ten T\ten R)=\tr\ten T$，并说明每一步用了什么性质。}

\textbf{先把用到的两个工具说清楚。}第一个是迹的定义：对 $3\times3$ 矩阵 $\ten A$，
\[
\tr\ten A=A_{11}+A_{22}+A_{33}=\sum_{i=1}^{3}A_{ii}.
\]
第二个是迹的循环性。题目已经把这条性质作为工具给出；这里把它写清楚，是为了让第一步有依据：
\[
\tr(\ten A\ten B)=\tr(\ten B\ten A)
\]
对任意两个同阶方阵 $\ten A,\ten B$ 成立。
证明是一次指标求和：
\[
\tr(\ten A\ten B)
=\sum_{i=1}^{3}(\ten A\ten B)_{ii}
=\sum_{i=1}^{3}\sum_{j=1}^{3}A_{ij}B_{ji}
=\sum_{j=1}^{3}\sum_{i=1}^{3}B_{ji}A_{ij}
=\sum_{j=1}^{3}(\ten B\ten A)_{jj}
=\tr(\ten B\ten A).
\]
这里只交换了两个有限和的次序，并对同一个乘积换了一种分组方式；没有用到对称、可逆或正交。$\ten A,\ten B$ 只要是同阶方阵即可。

\textbf{现在证明题目。}把 $\ten R^{\mathsf T}\ten T$ 看成一个整体，记作 $\ten A$，把 $\ten R$ 记作 $\ten B$。第一步用循环性：
\[
\tr\bigl((\ten R^{\mathsf T}\ten T)\ten R\bigr)
=\tr\bigl(\ten R(\ten R^{\mathsf T}\ten T)\bigr).
\]
这一步只用了 $\tr(\ten A\ten B)=\tr(\ten B\ten A)$，其中 $\ten A=\ten R^{\mathsf T}\ten T$，$\ten B=\ten R$。

第二步把右边的括号重新结合，并用正交性：
\[
\tr\bigl(\ten R(\ten R^{\mathsf T}\ten T)\bigr)
=\tr\bigl((\ten R\ten R^{\mathsf T})\ten T\bigr)
=\tr(\unitd\ten T)
=\tr\ten T .
\]
第一个等号用了矩阵乘法的结合律，把 $\ten R$ 与 $\ten R^{\mathsf T}$ 先乘起来；第二个等号用了正交性 $\ten R\ten R^{\mathsf T}=\unitd$；第三个等号用了单位矩阵 $\unitd\ten T=\ten T$。把两步接起来就是
\[
\tr(\ten R^{\mathsf T}\ten T\ten R)
=\tr\bigl((\ten R^{\mathsf T}\ten T)\ten R\bigr)
=\tr\bigl(\ten R(\ten R^{\mathsf T}\ten T)\bigr)
=\tr\bigl((\ten R\ten R^{\mathsf T})\ten T\bigr)
=\tr\ten T .
\]

\textbf{每一步用了什么性质。}第一步用“迹的循环性”；随后用“矩阵乘法的结合律”；再用“$\ten R$ 正交，即 $\ten R\ten R^{\mathsf T}=\unitd$”；最后用“单位矩阵乘任何矩阵不变”。题目只要求“写出两步”，核心就是“循环一次，再把 $\ten R\ten R^{\mathsf T}$ 结合出来”；中间多写的等号只是把结合律写细，不改变逻辑。

\textbf{条件与含义。}这里 $\ten R$ 必须是方阵且正交；$\ten T$ 不必对称，也不必可逆。结论说明：若 $\ten T'=\ten R^{\mathsf T}\ten T\ten R$ 是同一线性映射在另一组正交归一基下的矩阵，那么 $\tr\ten T'=\tr\ten T$；也就是说，\textbf{迹在正交换基下不变}。但要注意，它不说明 $\ten T$ 的每个矩阵元都不变——换基时矩阵元一般会变，只有迹这样的组合不变。当 $\ten T$ 又是实对称矩阵、$\ten R$ 取 §1.5 的主轴矩阵时，把 $\ten R^{\mathsf T}\ten T\ten R=\operatorname{diag}(\lambda_1,\lambda_2,\lambda_3)$ 代入，还能得到
\[
\tr\ten T=\lambda_1+\lambda_2+\lambda_3,
\]
这与 §1.3 用迹抽查特征值的做法一致。题目已经把迹的循环性作为工具给出；本题要求证明的是 $\tr(\ten R^{\mathsf T}\ten T\ten R)=\tr\ten T$，并说明每一步用了什么性质。

\end{document}

